\documentclass[10pt,a4paper]{amsart}
\usepackage[margin=19mm]{geometry}
\usepackage{amsmath,amssymb,mathtools}
\usepackage[scr=rsfs,cal=euler]{mathalfa}
\usepackage{xfrac}
\usepackage[unicode,hidelinks,bookmarks=false]{hyperref}
\newcommand{\ie}{i.e.\ }
\newcommand{\cf}{cf.\ }
\newcommand{\resp}{respectively\ }
\newcommand{\Subsection}{\subsection}
\providecommand{\nobreakdash}{\nobreak\hbox{-}}
\setlength{\emergencystretch}{2em}
\multlinegap0pt
\usepackage{dsfont}
\usepackage{amssymb}
\usepackage[shortlabels]{enumitem}
\usepackage{csquotes}
\usepackage{float}
\usepackage{mdframed}
\newtheorem{proposition}{Proposition}
\newtheorem{remark}{Remark}
\newtheorem{corollary}{Corollary}
\newtheorem{lemma}{Lemma}
\newtheorem{theorem}{Theorem}
\newcommand{\E}{\operatorname{E}}
\renewcommand{\P}{\operatorname{P}}
\newcommand{\PP}{\mathbb{P}}
\newcommand{\Q}{\operatorname{Q}}
\newcommand{\RR}{\mathbb{R}}
\newcommand{\U}{\operatorname{U}}
\newcommand{\eps}{\varepsilon}
\renewcommand{\k}{\kappa}
\newcommand{\mE}{\mathcal{E}}
\newcommand{\mH}{\mathcal{H}}
\newcommand{\mP}{\mathcal{P}}
\newcommand{\mT}{\mathcal{T}}
\newcommand{\ol}{\overline}
\newcommand{\sX}{\mathrm{X}}
\newcommand{\supp}{\operatorname{supp}}
\newcommand{\wh}{\widehat}
\title{A Minimax Theory of Nonparametric \\Regression Under Covariate Shift}
\author{Petr Zamolodtchikov}
\date{Source: arXiv:2603.05897v1 (CC BY 4.0)}
\begin{document}
\maketitle

\noindent\emph{Source and licence.} A Minimax Theory of Nonparametric \\Regression Under Covariate Shift. Version arXiv:2603.05897v1 (CC BY 4.0), distributed by arXiv under the Creative Commons Attribution 4.0 International licence, http://creativecommons.org/licenses/by/4.0/, as recorded in the arXiv licence metadata for this exact version and shown on the abstract page https://arxiv.org/abs/2603.05897v1. Full licence notice: \texttt{licenses/arxiv-2603.05897-CC-BY-4.0.txt}.\par\smallskip

\noindent\textit{Editorial scope.} Counted unit: the article's theorem th.rates.local.k.two.sample.proof (source0.tex lines 1538--1570). The article's complete original proof of it is delivered with it (source0.tex lines 1571--1832, one \texttt{proof} environment of 262 lines). No mathematics is written, repaired or re-derived: the statement and the proof are the article's own lines, in the article's own notation, and no retained line is substituted or re-flowed.  Retained uncounted as the source-local context the proof needs: lines 374--377; lines 934--944; lines 999--1011; lines 1012--1022; lines 1134--1137; lines 1139--1142; lines 1173--1180; lines 1226--1236; lines 1265--1272; lines 1293--1306; lines 1333--1358; lines 1507--1513.  Nothing was omitted from any retained span, so the package carries no omission record.  No retained line contains a citation, so the wrapper carries no bibliography and no bibliography entry.  No retained line contains a figure environment, an image-inclusion command or drawing code, so no image is delivered and the package declares no asset dependency.  Editorial formatting only, in addition: an AMS \texttt{amsart} preamble that restates the article's own declarations and macros used by the retained lines, namely the environments \texttt{proposition}, \texttt{remark}, \texttt{corollary}, \texttt{lemma}, \texttt{theorem} on separate counters, together with the standard packages those lines need.  The retained environments print in this document as Proposition 1, Remark 1, Corollary 1, Lemma 1, Theorem 1 in order of appearance, whereas the article prints the same items with the numbering of its own full document.  The complete native source of the article as distributed in the arXiv e-print for this exact version is bundled unchanged as \texttt{sources/195827/source0.tex}.
\begin{align}
    \label{eq.model}
    Y = f_*(X) + \eps,
\end{align}

\begin{proposition}
    \label{prop.neighbours.distance.bound.zeta}
    Let $\tau > 0.$ If $k \geq \lceil \log n\rceil, \ n \geq 2d,$ and $h_+ = h_+(k) := 2k/((1 - \lambda^+_\tau)n)$ with
    \begin{align*}
        \lambda_\tau^+ := \frac{4\big(2d + \tau + \log(K_0)/\log(2d))\big) + 9}{4\big(2d + \tau + \log(K_0)/\log(2d))\big) + 12},
    \end{align*} 
    then,
    \begin{align*}
        \PP\Bigg\{\sup_{x \in \RR^d}\frac{R_k(x)}{\zeta_{h_+}(x)} \geq 1\Bigg\} \leq n^{-\tau}.
    \end{align*}
\end{proposition}

\begin{remark}
    \label{rem.bound.variable.k}
    While Proposition \ref{prop.neighbours.distance.bound.zeta} shows that $R_k(x) \leq \zeta_{h_+}(x)$ uniformly with high probability for each individual choice of $k \geq \log n,$ it can easily be extended to the case where $k = k(x) \in \{1, \dots, n\}$ depends on $x \in \RR^d.$ The fact that $k$ can now change values depending on $x$ demands one to pay the price that the $\sup$ in the probability is not on the whole $\RR^d$ anymore, but only on a subset $\Theta_n^*(k)$ defined as follows
    \begin{align*}
        \Theta_n(k) := \bigg\{x \in \RR^d: k(x) \geq \lceil\log n \rceil\bigg\}.
    \end{align*}
    Similarly, we have to pay the price of a factor $n$ in the probability, as can be seen from the union bound
    \begin{align*}
        \PP\Bigg\{\sup_{x \in \Theta_n^*}\frac{R_{k(x)}(x)}{\zeta_{h_+}} \geq 1\Bigg\} &\leq \PP\Bigg\{\bigcup_{\ell = \lceil \log n \rceil}^n\Bigg\{\sup_{x \in \RR^d}\frac{R_{\ell}(x)}{\zeta_{h_+}(x)} \geq 1\Bigg\} \Bigg\}\\
        &\leq \sum_{\ell = \lceil \log n \rceil}^n\PP\Bigg\{\sup_{x \in \RR^d}\frac{R_{\ell}(x)}{\zeta_{h_+}(x)}\geq 1\Bigg\}\\
        &\leq n^{1 - \tau}.
    \end{align*}
\end{remark}

\begin{proposition}
    \label{prop.neighbours.distance.lower.bound.zeta}
    Let $\tau > 0$ and $V_\tau := K_0^{1/(2d)}\vee(8d + 4\tau).$ If $k \geq \lceil V_\tau\log n\rceil, \ n\geq 2d,$ and $h_- = h_-(k) := k/(2n(1 + \lambda_\tau^-)),$ with 
    \begin{align*}
        \lambda_\tau^- := 1 + \sqrt{\frac 52},
    \end{align*}
    then,
    \begin{align*}
        \PP\Bigg\{\inf_{x \in \RR^d}\frac{R_k(x)}{\zeta_{h_-}(x)} \leq 1\Bigg\} \leq n^{-\tau}.
    \end{align*}
\end{proposition}

\begin{align}
    \label{eq.local.knn}
    \wh f(x) = \frac{1}{k(x)}\sum_{i=1}^n Y_i(x),
\end{align}

\begin{align}
    \label{eq.average.knn}
    \ol f\colon x \mapsto \E\big[\wh f(x)|X_1, \dots, X_n\big].
\end{align}

\begin{proposition}\label{prop.variance.bound.general.case}
Let $\wh f$ be as in \eqref{eq.local.knn}, with neighbour function $k \colon \RR^d \to \{1, \dots, n\}.$ Let $\ol f$ be defined by \eqref{eq.average.knn}.
Then, for all $n\geq 2d,$ it holds that
\begin{align*}
    \PP\Bigg\{\sup_{x \in \RR^d}k(x)^{1/2}\big|\wh f(x) - \ol f(x)\big| \geq \eta\Bigg\} &\leq K_dn^{2d + 1}e^{-H \eta^2},
\end{align*}
where $K_d = 2(25/d)^d,$ and $H = \lambda_0^2/(8\sigma_0)$ for arbitrary $\lambda_0, \sigma_0 > 0$ such that $\E[\exp(\lambda|\eps|)] \leq \sigma_0.$
\end{proposition}

\begin{corollary}
    \label{cor.variance.bound.nice.version}
    Under the conditions of Proposition \ref{prop.variance.bound.general.case}, we have for all $\tau > 0,$ all $n \geq 3 \vee 2d,$
    \begin{align*}
        \PP\Bigg\{\sup_{x \in \RR^d}k(x)^{1/2}\big|\wh f(x) - \ol f(x)\big|\geq \sqrt{C_\tau \log n}\Bigg\} \leq n^{-\tau},
    \end{align*}
    with 
    \begin{align*}
        C_\tau := \frac{\log K_d}{H}\bigg(1 + \frac{2d + \tau + 1 }{\log(K_d))}\bigg)
    \end{align*} where $H$ and $K_d$ are defined in Proposition \ref{prop.variance.bound.general.case}.
\end{corollary}

\begin{proposition}[Bias-Variance decomposition]
    \label{prop.decomp.knn}
    Under the assumptions of model \eqref{eq.model}, if $\wh f$ is a local $k$-NN with neighbour function $k\colon \RR^d \to \{\lceil \log n\rceil, \dots, n\},$ and if $\P \in \mP,$ then, for all $n \geq 2d \vee 3$ and all $\Delta \subseteq \Theta_n^{\P}(S_\tau, k),$
    \begin{multline*}
        \Bigg\{\forall x \in \RR^d,\ \big(\wh f(x) - f(x)\big)^2 \leq 2C_\tau\frac{\log n}{k(x)} + 8L^2\bigg(\frac{S_\tau k(x)}{np(x)}\bigg)^{2\beta/d}\mathds{1}(x \in \Delta) + 8L^2\mathds{1}(x \in \Delta^c)\Bigg\}\\
        \supseteq \bigg(\bigcap_{i \in k(\RR^d)}U_n^{\P}(i)\bigg)\bigcap V_n^{\P}(\wh f).
    \end{multline*}
\end{proposition}

\begin{lemma}
    \label{lem.density.ratio}
    Let $\P \in \mP(D, \theta)$ with $D, \theta > 0.$ Denote by $p$ the density of $\P$, and let $\wh p$ be the $\ell$-NN density estimator obtained from $n$ i.i.d.\ samples from $\P.$ Work with $n \geq 3\vee 2d.$
    \begin{enumerate}
        \item [$(i)$] If $\ell \geq \lceil V_\tau\log n\rceil,$ then
        \begin{align*}
            L_n^{\P}(\ell) \subseteq \Bigg\{\sup_{x \in \Delta_n^{\P}(c_2, \ell)} \frac{\wh p(x)}{p(x)} \leq \frac{\theta}{c_1}\Bigg\}.
        \end{align*}
        \item [$(ii)$] If $\ell \geq \lceil \log n \rceil,$ then
        \begin{align*}
            U_n^{\P}(\ell) \subseteq \Bigg\{\inf_{x \in \Delta_n^{\P}(S_\tau, \ell)} \frac{\wh p(x)}{p(x)} \geq \frac{1}{S_\tau}\Bigg\}
        \end{align*}
    \end{enumerate}
\end{lemma}

\begin{lemma}
    \label{lem.inclusions}
    Let $\P \in \mP,\ \tau > 0$ and $\k > 0.$ Consider the function
    \begin{align*}
        k(x) = n \wedge \Big\lceil\k\log(n)^{d/(2\beta + d)}\big(n\wh p(x)\big)^{2\beta/(2\beta + d)} \Big\rceil \vee \lceil \log n\rceil, 
    \end{align*}
    where $\wh p$ is a $\ell$-NN density estimator for $\P,$ with $\ell \geq \lceil V_\tau\log n\rceil.$ Let $N = N(\k, \theta, D) \geq 3\vee 2d$ satisfy
    \begin{align*}
        \frac N{\log N} \geq \k^{(2\beta + d)/d}\bigg(\frac{\theta D}{c_1}\bigg)^{2\beta/d},
    \end{align*}
    and define the constant
    \begin{align*}
        C := S_\tau\vee c_2 \vee \frac{S_\tau}{\k^{(2\beta + d)/(2\beta)}V_\tau} \vee \frac{(2S_\tau\k)^{(2\beta + d)/d}}{V_\tau}\bigg(\frac{\theta}{c_1}\bigg)^{2\beta/d}.
    \end{align*}
    Then, the following assertions hold
    \begin{enumerate}
        \item [$(i)$] For all $n \geq N,$ 
            \begin{align*}
                L_n^{\P}(\ell)\cap U_n^{\P}(\ell) \subseteq \Bigg\{\forall x \in \Delta_n^{\P}(C, \ell), \ k(x) = \Big\lceil\k\log(n)^{d/(2\beta + d)}\big(n\wh p(x)\big)^{2\beta/(2\beta + d)}\Big\rceil\Bigg\}.
            \end{align*}
        \item [$(ii)$] For all $n \geq N,$
            \begin{align*}
                 L_n^{\P}(\ell)\cap U_n^{\P}(\ell) \subseteq \Bigg\{\Delta_n^{\P}(C, \ell) \subseteq \Theta_n^{\P}(S_\tau, k)\Bigg\}.
            \end{align*}
    \end{enumerate}
\end{lemma}

\begin{lemma}
    \label{lem.bounded.below.zeta.low.density}
    Let $\P \in \mP(D, \theta)$ with $D, \theta > 0,$ and $k \in \{1, \dots, n\}.$ If $x \in \supp(\P)$ satisfies $p(x) \leq \alpha k/n,$ with $\alpha > 0,$ then, 
    \begin{align*}
         \zeta_{h_-}(x) \geq \bigg(\frac{c_1}{\alpha \theta}\bigg)^{1/d} \wedge 1.
    \end{align*}
\end{lemma}

\begin{theorem}[Rates for two-sample local nearest neighbours]
\label{th.rates.local.k.two.sample.proof}
    Let $\tau > 0, \k_{\P}, \k_{\Q} >0,$ and define
    \begin{align*}
        r_S := 2\beta/(2\beta + d) \wedge \gamma,\text{ and }r_T := s\wedge 2\beta/(2\beta + d).
    \end{align*}
    Let $\wh p$ and $\wh q$ be, respectively, the $\ell_{\P}$ and $\ell_{\Q}$-NN density estimators of $p$ and $q$ with
    \begin{align*}
        \ell := \ell_{\P} = \ell_{\Q} = \big\lceil V_\tau\log nm\big\rceil.
    \end{align*}
    Finally, consider $\wh f$ to be the local $(k_{\P}, k_{\Q})$-NN estimator with neighbour functions given by
    \begin{align*}
        k_{\P}(x) &:= n\wedge \Big\lceil \log(nm)^{d/(2\beta + d)}(n\wh p(x))^{2\beta/(2\beta + d)}\Big\rceil \vee \big\lceil \log (nm) \big\rceil\\
        k_{\Q}(x) &:= m\wedge \Big\lceil \log(nm)^{d/(2\beta + d)}(m\wh q(x))^{2\beta/(2\beta + d)}\Big\rceil \vee \big\lceil \log (nm) \big\rceil.
    \end{align*}
    Then, there exists $N \geq 2d$ and $M \geq 2d$ such that for all $n \geq N, m \geq M,$ it holds, with probability at least $1 - 3n^{1-\tau} - 3m^{1-\tau},$
    \begin{align*}
        \sup_{f \in \mH(L, \beta)}\PP\Big\{\mE_{\Q}(\wh f, f) \geq C(\tau)R(n,m, \gamma, s)\Big\} \leq 3(n^{1 - \tau} + m^{1 - \tau})
    \end{align*}
    with, if $(\gamma - r_\beta)(s - r_\beta) <0$ and $m \in [n, n^{\gamma/s}],$ then
    \begin{align*}
        R(n,m, \gamma, s) &= \tilde C_0\Big(\frac{\log(nm)}{n}\Big)^{\gamma a}\Big(\frac{\log(nm)}{m}\Big)^{s(1- a)}\mT_{\P}(\gamma)^{a}\mT_{\Q}(s)^{1-a},
    \end{align*}
    where $\tilde C_0$ is a constant independent of $n,m, \gamma, s, \tau$ and, if $(\gamma - r_\beta)(s - r_\beta) \geq 0$ or $m \notin [n, n^{\gamma/s}],$ then
    \begin{align*}
         R &= \ol C_0\Big[\mT_{\P}(r_S)\Big(\frac{\log(nm)}{n}\Big)^{r_S}\Big]\wedge \Big[\mT_{\Q}(r_T)\Big(\frac{\log(nm)}{m}\Big)^{r_T}\Big],
    \end{align*}
    where $\ol C_0$ is a constant independent of $n,m, \gamma, s, \tau$ and $C_p := C(\k_{\P}), \ C_q := C(\k_{\Q}),$ where
    \begin{align*}
        C(t) := S_\tau\vee c_2 \vee \frac{S_\tau}{t^{(2\beta + d)/(2\beta)}V_\tau} \vee \frac{(2S_\tau t)^{(2\beta + d)/d}}{V_\tau}\bigg(\frac{\theta}{c_1}\bigg)^{2\beta/d}.
    \end{align*}
    
\end{theorem}

\begin{proof}
    Let 
    \begin{align*}
        C(t) := S_\tau\vee c_2 \vee \frac{S_\tau}{t^{(2\beta + d)/(2\beta)}V_\tau} \vee \frac{(2S_\tau t)^{(2\beta + d)/d}}{V_\tau}\bigg(\frac{\theta}{c_1}\bigg)^{2\beta/d}.
    \end{align*}
    and define $C_p := C(\k_{\P}), \ C_q := C(\k_{\Q}).$ For $R \in \{\P, \Q\},$ Define the event
    \begin{align*}
        E_n^{R_{\sX}}(\ell, \wh f_R) := L_n^{R_{\sX}} \cap \U_n^{R_{\sX}}(\ell) \cap V_n^{R_{\sX}}(\wh f_R).
    \end{align*}
    We work on the event
    \begin{align}
        \label{eq.event.two.sample.local}
        E_n^{\P_{\sX}}(\ell_{\P}, \wh f_{\P}) \cap E_m^{\Q\mathstrut_{\!\sX}}(\ell_{\Q}, \wh f_{\Q}).
    \end{align}
   Let $X \sim \Q\mathstrut_{\!\sX}$ be independent of the samples. We start by decomposing the risk as follows.
    \begin{align*}
        \big\|\wh f - f_*\big\|_{L^2(\Q\mathstrut_{\!\sX})}^2 &= \E_X\bigg[\bigg(\frac{1}{k_{\P}(X) + k_{\Q}(X)}\bigg(\sum_{i=1}^{k_{\P}(X)}\big(Y_i(X) - f_*(X)\big) + \sum_{j=1}^{k_{\Q}(X)} \big(Y_j'(X) - f_*(X)\big)\bigg)\bigg)^2\bigg]\\
        &= \E\bigg[\bigg(\frac{k_{\P}(X)}{k_{\P}(X) + k_{\Q}(X)}\big(\wh f_{\P}(X) - f_*(X)\big) + \frac{k_{\Q}}{k_{\P}(X) + k_{\Q}(X)}\big(\wh f_{\Q}(X) - f_*(X)\big)\bigg)^2\bigg],
    \end{align*}
    where $\wh f_{\P}(x) := k_{\P}(x)^{-1}\sum_{i=1}^{k_{\P}(x)} Y_i(x)$ and $\wh f_{\Q}(x) := k_{\Q}(x)^{-1}\sum_{j=1}^{k_{\Q}(x)} Y_j'(x).$ Denote by $w_{\P}(x) := k_{\P}(x)/(k_{\P}(x) + k_{\Q}(x))$ and by $w_{\Q}(x) := k_{\Q}(x)/(k_{\P}(x) + k_{\Q}(x)).$ A further application of Jensen's inequality yields
    \begin{align}
        \label{eq.decomposition.expectation.two.sample}
        \big\|\wh f - f_*\big\|_{L^2(\Q\mathstrut_{\!\sX})}^2 &\leq \E_X\Big[w_{\P}(X)\big(\wh f_{\P}(X) - f_*(X)\big)^2 + w_{\Q}(X) \big(\wh f_{\Q}(X) - f_*(X)\big)^2\Big].
    \end{align}
    Let $\Delta_{\P} := \Delta_n^{\P_{\sX}}(C_p, \ell_{\P})$ and $\Delta_{\Q} := \Delta_m^{\Q\mathstrut_{\!\sX}}(C_q, \ell_{\Q}).$ We apply Proposition \ref{prop.decomp.knn} twice to obtain bounds on $(\wh f_{\P}(X) - f_*(X))^2$ and $(\wh f_{\Q}(X) - f_*(X))^2.$ On the event \eqref{eq.event.two.sample.local}, we can decompose the error inside the expectation in \eqref{eq.decomposition.expectation.two.sample} in four different
    ways, depending on the region. The result is as follows
    \begin{align}
        \label{eq.first.bv}
        \begin{split}
        \big(\wh f_{\P}(x) - f_*(x)\big)^2 &\leq \begin{dcases}
            b^2_{\P}(x) + V_{\P}(x) & \text{ if } x \in \Delta_{\P}\\
            8L^2 + V_{\P}(x) &\text{ if } x \in \Delta_{\P}^c
        \end{dcases}\\
        \big(\wh f_{\Q}(x) - f_*(x)\big)^2 &\leq \begin{dcases}
            b^2_{\Q}(x) + V_{\Q}(x) & \text{ if } x \in \Delta_{\Q}\\
            8L^2 + V_{\Q}(x) &\text{ if } x \in \Delta_{\Q}^c,
        \end{dcases}
        \end{split}
    \end{align}
    where
    \begin{align*}
        b_{\P}^2(x) &:= 8L^2\bigg(\frac{S_\tau k_{\P}(x)}{np(x)}\bigg)^{2\beta/d}, \ V_{\P}(x) := 2C_\tau \frac{\log n}{k_{\P}(x)},\\
        b_{\Q}^2(x) &:= 8L^2\bigg(\frac{S_\tau k_{\Q}(x)}{mq(x)}\bigg)^{2\beta/d}, \ V_{\Q}(x) := 2C_\tau \frac{\log m}{k_{\Q}(x)},
    \end{align*}
    In order to derive bounds on the excess risk, we further subdivide $\RR^d$ into four regions. Let 
    \begin{align*}
        \Delta := \Delta_{\P}\cap \Delta_{\Q}, \ \Gamma_{\P} := \Delta_{\P}\cap \Delta_{\Q}^c,\ \Gamma_{\Q} := \Delta_{\P}^c\cap \Delta_{\Q},\text{ and } \Theta := \Delta_{\P}^c\cap \Delta_{\Q}^c.
    \end{align*}
    We now bound the pointwise risk of $\wh f$ on each of these regions.\\
    \noindent\textbf{Region $\Delta$.} For $x \in \Delta,$ the total pointwise risk is
    \begin{align}
        \label{eq.pointwise.loss.bulk}
        \big(\wh f(x) - f_*(x)\big)^2 &\leq w_{\P}(x)\big(b_{\P}(x)
        ^2 + V_{\P}(x)\big) + w_{\Q}(x)\big(b_{\Q}(x)
        ^2 + V_{\Q}(x)\big).
    \end{align}
    We now denote $r_\beta := 2\beta/(2\beta + d), \ \mT_{\P}(u) := \int q(x)/p(x)^{u}\, dx$ and $\mT_{\Q}(v) := \int q(x)^{1 -v}\, dx.$ From the definitions of $k_{\P}, k_{\Q},$ Lemma \ref{lem.inclusions} and Lemma \ref{lem.density.ratio} $(ii)$, we obtain
    \begin{align}
        \label{eq.two.sample.local.variance.bound}
        w_{\P}(x)V_{\P}(x) + w_{\Q}(x)V_{\Q}(x) &= 2C_\tau\Big(\frac{\log n}{k_{\P}(x) + k_{\Q}(x)} + \frac{\log m}{k_{\P}(x) + k_{\Q}(x)}\Big)\nonumber\\
        &= 2C_\tau\frac{\log (nm)}{k_{\P}(x) + k_{\Q}(x)}\\
        & = 2C_\tau \Big(\frac{k_{\P}(x)}{\log (nm)} + \frac{k_{\Q}(x)}{\log (nm)}\Big)^{-1}\nonumber\\
        &= 2C_\tau \Bigg(\k_{\P}\frac{\log(nm)^{1 - r_\beta}(n\wh p(x))^{r_\beta}}{\log(nm)} + \k_{\Q}\frac{\log(nm)^{1 - r_\beta}(m\wh q(x))^{r_\beta}}{\log(nm)}\Bigg)^{-1}\nonumber\\
        &\leq 2C_\tau \Bigg(\k_{\P}\bigg(\frac{np(x)}{S_\tau\log(nm)}\bigg)^{r_\beta} + \k_{\Q}\bigg(\frac{m q(x)}{S_\tau \log(nm)}\bigg)^{r_\beta}\Bigg)^{-1}\nonumber\\
        &= 2C_\tau S_\tau^{r_\beta} \Bigg(\k_{\P}\bigg(\frac{np(x)}{\log(nm)}\bigg)^{r_\beta} + \k_{\Q}\bigg(\frac{m q(x)}{\log(nm)}\bigg)^{r_\beta}\Bigg)^{-1}
    \end{align}
    We now bound the remaining terms in \eqref{eq.pointwise.loss.bulk}. From the definition of $k_{\P},$ and Lemma \ref{lem.density.ratio}, we get
    \begin{align*}
        k_{\P}(x)b_{\P}^2(x) = k_{\P}(x)\bigg(\frac{S_\tau k_{\P}(x)}{np(x)}\bigg)^{2\beta/d} &= S_\tau^{2\beta/d}\frac{k_{\P}(x)^{(2\beta + d)/d}}{(np(x))^{2\beta/d}}\\
        &=S_\tau^{2\beta/d}\k_{\P}^{(2\beta + d)/d}\log(nm)\bigg(\frac{\wh p(x)}{p(x)}\bigg)^{2\beta/d}\\
        &\leq \k_{\P}^{(2\beta + d)/d}\bigg(\frac{\theta S_\tau}{c_1}\bigg)^{2\beta/d}\log(nm) := K_{\P}\log(nm).
    \end{align*}
    The same analysis for the second term in the expression of $b^2(x)$ yields
    \begin{align*}
        k_{\Q}(x)b_{\Q}^2(x) = k_{\Q}(x)\bigg(\frac{S_\tau k_{\Q}(x)}{mq(x)}\bigg)^{2\beta/d} \leq \k_{\Q}^{(2\beta + d)/d}\bigg(\frac{\theta S_\tau}{c_1}\bigg)^{2\beta/d}\log(nm) =: K_{\Q}\log(nm).
    \end{align*}
    Hence,
    \begin{align*}
        8L^2(w_{\P}(x)b^2_{\P}(x) + w_{\Q}(x)b^2_{\Q}(x)) &\leq 8L^2(K_{\P} + K_{\Q})\frac{\log(nm)}{k_{\P}(x) + k_{\Q}(x)} =: \ol K\frac{\log(nm)}{k_{\P}(x) + k_{\Q}(x)},
    \end{align*}
    where $\ol K := K_{\P} + K_{\Q}.$ The latter quantity can be bounded in the exact same way as in \eqref{eq.two.sample.local.variance.bound}, which shows that, for $x \in \Delta,$
    \begin{align*}
        \big(\wh f(x) - f_*(x)\big)^2 &\leq (2C_\tau + 8L^2\ol K)S_\tau^{r_\beta} \Bigg(\k_{\P}\bigg(\frac{np(x)}{\log(nm)}\bigg)^{r_\beta} + \k_{\Q}\bigg(\frac{m q(x)}{\log(nm)}\bigg)^{r_\beta}\Bigg)^{-1}\\
        &:= C_{\Delta}\Bigg(\k_{\P}\bigg(\frac{np(x)}{\log(nm)}\bigg)^{r_\beta} + \k_{\Q}\bigg(\frac{m q(x)}{\log(nm)}\bigg)^{r_\beta}\Bigg)^{-1}\\
        &\leq C_{\Delta}\bigg(\k_{\P}^{-1}\Big(\frac{\log(nm)}{np(x)}\Big)^{r_\beta}\wedge\k_{\Q}^{-1}\Big(\frac{\log(nm)}{mq(x)}\Big)^{r_\beta}\bigg)
    \end{align*}
    \noindent\textbf{Regions $\Gamma_{\P}$ and $\Gamma_{\Q}.$} For $x \in \Gamma_{\P},$ we have
    \begin{align*}
        \big(\wh f(x) - f_*(x)\big)^2 &\leq w_{\P}\big(b_{\P}^2(x) + V_{\P}(x)\big)^2 + w_{\Q}(x)\big(8L^2 + V_{\Q}(x)\big).
    \end{align*}
    Since $x \in \Delta_{\P},$ the bound on $w_{\P}(x)(\wh f_{\P}(x) - f_*(x))^2$ from the previous case still holds, and we obtain
    \begin{align*}
        \big(\wh f(x) - f_*(x)\big)^2 &\leq \frac{(K_{\P} + 2C_\tau) \log(nm) + 8L^2k_{\Q}(x)}{k_{\P}(x) + k_{\Q}(x)}.
    \end{align*}
    Recall that $x \in \Delta_{\Q}^c,$ meaning that $q(x) \leq C_q\ell_{\Q}/m,$ hence, by virtue of Lemma \ref{lem.bounded.below.zeta.low.density}, it holds that
    \begin{align}
        \label{eq.ub.wh.q.low.density}
        \wh q(x) = \frac{\ell_{\Q}}{mR_{\ell_{\Q}, \Q}(x)^d} \leq \frac{\ell_{\Q}}{m\zeta_{h_-, \Q}(x)^d} \leq \frac{a\theta}{c_1}\frac{\ell_{\Q}}{m}.
    \end{align}
    Hence, recalling that $\ell_{\P} = \ell_{\Q} = \lceil V_\tau \log(nm)\rceil,$
    \begin{align*}
        \big(\wh f(x) - f_*(x)\big)^2 &\leq \frac{(K_{\P} + 2C_\tau) \log(nm) + 16L^2\k_{\Q}(a\theta/c_1)^{r_\beta}\log(nm)^{1-r_\beta}\ell_{\Q}^{r_\beta}}{k_{\P}(x) + k_{\Q}(x)}\\
        &\leq (K_{\P} + 2C_\tau + 16L^2V_\tau^{r_\beta}\k_{\Q}(a\theta/c_1)^{r_\beta}) \frac{\log(nm)}{k_{\P}(x) + k_{\Q}(x)}\\
        &\leq (K_{\P} + 2C_\tau + 16L^2V_\tau^{r_\beta}\k_{\Q}(a\theta/c_1)^{r_\beta})\frac{\log(nm)}{k_{\P}(x)}\\
        &\leq \frac{K_{\P} + 2C_\tau + 16L^2V_\tau^{r_\beta}\k_{\Q}(a\theta/c_1)^{r_\beta}}{\k_{\P}}\bigg(\frac{\log(nm)}{n\wh p(x)}\bigg)^{r_\beta}\\
        &\leq \frac{(K_{\P} + 2C_\tau + 16L^2V_\tau^{r_\beta}\k_{\Q}(a\theta/c_1)^{r_\beta})S_\tau^{r_\beta}}{\k_{\P}}\bigg(\frac{\log(nm)}{n p(x)}\bigg)^{r_\beta}=: C_{\Gamma_{\P}}\bigg(\frac{\log(nm)}{np(x)}\bigg)^{r_\beta}.
    \end{align*}
    The analysis on the region $\Gamma_{\Q}$ is symmetric and we obtain, for $x \in \Gamma_{\Q},$
    \begin{align*}
        \big(\wh f(x) - f_*(x)\big)^2 &\leq C_{\Gamma_{\Q}}\bigg(\frac{\log(nm)}{mq(x)}\bigg)^{r_\beta}, \ \text{ with } C_{\Gamma_{\Q}} := \frac{(K_{\Q} + 2C_\tau + 16L^2V_\tau^{r_\beta}\k_{\P}(a\theta/c_1)^{r_\beta})S_\tau^{r_\beta}}{\k_{\Q}}.
    \end{align*}
    \noindent\textbf{Region $\Theta.$} For $x \in \Theta,$ we use the fact that, by definition, $k_{\P}(x)\wedge k_{\Q}(x) \geq \log(nm)$ to obtain
    \begin{align*}
        \big(\wh f(x) - f_*(x)\big)^2 &\leq 8L^2 + w_{\P}(x) V_{\P}(x) + w_{\Q}(x)V_{\Q}(x)\\
        &\leq 8L^2 + 2C_\tau \frac{\log (nm)}{k_{\P}(x) + k_{\Q}(x)}\\
        &\leq 8L^2 + C_\tau =: C_{\theta}
    \end{align*}
    This finishes the pointwise bounds on our partition of $\RR^d.$ It remains to integrate the latter against $\Q\mathstrut_{\!\sX}.$ We will proceed to integrate in each region. For each region, we will distinguish cases between the \textbf{wedge regime} and the \textbf{accelerated regime}. We define
    \begin{align*}
        I_{\Delta} &:= \int_{\Delta}\bigg(\k_{\P}^{-1}\Big(\frac{1}{np(x)}\Big)^{r_\beta}\wedge\k_{\Q}^{-1}\Big(\frac{1}{mq(x)}\Big)^{r_\beta}\bigg)\, d\Q\mathstrut_{\!\sX}(x)\\
        I_{\Gamma_{\P}} &:= \int_{\Gamma_{\P}}\bigg(\frac{1}{np(x)}\bigg)^{r_\beta}\, d\Q\mathstrut_{\!\sX}\\
        I_{\Gamma_{\Q}} &:= \int_{\Gamma_{\Q}}\bigg(\frac{1}{mq(x)}\bigg)^{r_\beta}\, d\Q\mathstrut_{\!\sX}(x)\\
        I_{\Theta} &:= \int_{\Theta}\, d\Q\mathstrut_{\!\sX}(x).
    \end{align*}
    With these notations, the final bound on the excess risk reads
    \begin{align*}
        \big\|\wh f - f_*\big\|_{L^2(\Q\mathstrut_{\!\sX})}^2 &\leq \log(nm)^{r_\beta}\big(C_{\Delta}I_{\Delta} + C_{\Gamma_{\P}}I_{\Gamma_{\P}} + C_{\Gamma_{\Q}}I_{\Gamma_{\Q}}) + C_{\Theta}I_{\Theta}.
    \end{align*}
    
    \noindent\textbf{Bound on $\Delta.$}
    We write
    \begin{align*}
        I_{\Delta} &\leq \inf_{\theta \in [0, 1]}\int_{\Delta}(\k_{\P}(np(x))^{r_\beta})^{-\theta}(\k_{\Q}(mq(x))^{r_\beta})^{1 - \theta}\, d\Q\mathstrut_{\!\sX}(x)\\
        &\leq \inf_{\theta \in [0, 1]}(\k_{\P}n^{r_\beta})^{-\theta}(\k_{\Q}m^{r_\beta})^{-(1 - \theta)}\int_{\Delta}p(x)^{-r_\beta \theta}q(x)^{-r_\beta(1 - \theta)}\, d\Q\mathstrut_{\!\sX}\\
        &\leq \inf_{\theta \in [0, 1]}(\k_{\P}n^{r_\beta})^{-\theta}(\k_{\Q}m^{r_\beta})^{-(1 - \theta)}\E_{X \sim \Q\mathstrut_{\!\sX}}\big[p(X)^{-r_\beta \theta}q(X)^{-r_\beta(1 - \theta)}\big]
    \end{align*}
    In the \textbf{wedge regime,} it is sufficient to use the facts that $r_S - r_\beta \leq 0, \ r_M - r_\beta \leq 0$ and that, on $\Delta, \ p(x) \geq C_p\ell_{\P}/n$ and $q(x) \geq C_q \ell_{\Q}/m$ and 
    \begin{align*}
        I_{\Delta} &\leq \bigg[\frac{1}{\k_{\P}n^{r_\beta}}\int_{\Delta}\frac{q(x)}{p(x)^{r_\beta}}\, dx\bigg]\wedge \bigg[\frac{1}{\k_{\Q}m^{r_\beta}}\int_{\Delta}q(x)^{1 - r_\beta}\, dx\bigg]\\
        &\leq \bigg[\frac{1}{\k_{\P}n^{r_\beta}}\int_{\Delta}\frac{q(x)}{p(x)^{r_S}}p(x)^{r_S - r_\beta}\, dx\bigg]\wedge \bigg[\frac{1}{\k_{\Q}m^{r_\beta}}\int_{\Delta}q(x)^{1 - r_\beta}\, dx\bigg]\\
        &\leq \bigg[\frac{1}{\k_{\P}n^{r_S}}(C_p\ell_{\P})^{r_S - r_\beta}\mT_{\P}(r_S)\bigg]\wedge\bigg[\frac{1}{\k_{\Q}m^{r_T}}(C_q\ell_{\Q})^{r_T - r_\beta}\mT_{\Q}(r_T)\bigg].
    \end{align*}
    In the \textbf{acceleration regime}, we can instead take $\theta = \gamma(s - r)/(r_\beta(s - \gamma))$ to obtain
    \begin{align*}
        I_{\Delta} &\leq (\k_{\P}n^{r_\beta})^{-\theta}(\k_{\Q}m^{r_\beta})^{-(1 - \theta)}\E_{X \sim \Q\mathstrut_{\!\sX}}\big[p(X)^{-r_\beta \theta}q(X)^{-r_\beta(1 - \theta)}\big]\\
        &\leq \k_{\P}^{-\theta}\k_{\Q}^{-(1 - \theta)}n^{-\tfrac{\gamma(r_\beta - s)}{\gamma - s}}m^{-\tfrac{s(\gamma - r)}{\gamma - s}}\E_{X \sim \Q\mathstrut_{\!\sX}}\big[p(X)^{-r_\beta \theta}q(X)^{-r_\beta(1 - \theta)}\big].
    \end{align*}
    It remains to verify that the expectation above is finite. We apply H\"older's inequality with the conjugates $a := \gamma/(\theta r_\beta),$ and $b := s/((1 - \theta)r_\beta),$ which leads to
    \begin{align*}
        \E_{X \sim \Q\mathstrut_{\!\sX}}\big[p(X)^{-r_\beta \theta}q(X)^{-r_\beta(1 - \theta)}\big] &\leq \Big(\E_{X \sim \Q\mathstrut_{\!\sX}}\big[p(X)^{-r_\beta \theta a}\big]\Big)^{1/a}\Big(\E_{X \sim \Q\mathstrut_{\!\sX}}\big[q(X)^{-r_\beta(1 - \theta)b}\big]\Big)^{1/b},
    \end{align*}
    and, since $r_\beta\theta a = \gamma$ and $r_\beta(1 - \theta)b = s,$ we conclude that
    \begin{align*}
        I_{\Delta} &\leq \k_{\P}^{-\theta}\k_{\Q}^{-(1 - \theta)}n^{-\tfrac{\gamma(r_\beta - s)}{\gamma - s}}m^{-\tfrac{s(\gamma - r)}{\gamma - s}}\mT_{\P}(\gamma)^{\theta r_\beta/\gamma} \mT_{\Q}(s)^{(1 - \theta)r_\beta/s}.
    \end{align*}
    
    \noindent\textbf{Bound on $\Gamma_{\P}.$} On $\Gamma_{\P},$ by definition, we have
    \begin{align*}
        np(x) \geq C_p\ell_{\P}.
    \end{align*}
    Therefore, in the \textbf{wedge regime}, it holds that
    \begin{align*}
        I_{\Gamma_{\P}} &= \int_{\Gamma_{\P}}(np(x))^{-r_\beta}\, d\Q\mathstrut_{\!\sX} \leq \int_{\Gamma_{\P}}\big[(C_p\ell_{\P})^{-r_\beta} \wedge (np(x))^{-r_\beta}\big]\, d\Q\mathstrut_{\!\sX}.
    \end{align*}
    Since for $x \in \Gamma_{\P}, \ q(x) \leq C_q\ell_{\Q}/m,$ and $p(x) \geq C_p\ell_{\P}/n,$ it holds that 
    \begin{align*}
        I_{\Gamma_{\P}} &\leq \bigg[\frac{(C_p\ell_{\P})^{r_S - r_\beta}}{n^{r_S}}\int_{\Gamma_{\P}}p(x)^{-r_S}\, d\Q\mathstrut_{\!\sX}(x)\bigg] \wedge \bigg[(C_p\ell_{\P})^{-r_\beta}\int_{\Gamma_{\P}}\frac{q(x)}{q(x)^{r_T}}q(x)^{r_T}\, dx\bigg]\\
        &\leq \bigg[\frac{(C_p\ell_{\P})^{r_S - r_\beta}}{n^{r_S}}\mT_{\P}(r_S)\bigg] \wedge \bigg[(C_p\ell_{\P})^{-r_\beta}\int_{\Gamma_{\P}}\frac{q(x)}{q(x)^{r_T}}q(x)^{r_T}\, dx\bigg]\\
        &\leq \bigg[\frac{(C_p\ell_{\P})^{r_S - r_\beta}}{n^{r_S}}\mT_{\P}(r_S)\bigg] \wedge \bigg[(C_p\ell_{\P})^{-r_\beta}\Big(\frac{C_q\ell_{\Q}}{m}\Big)^{r_T}\mT_{\Q}(r_T)\bigg],
    \end{align*}
    which gives the bound in the \textbf{wedge regime}. Next, in the \textbf{accelerated regime} where either $s < r_\beta < \gamma,$ or $\gamma < r_\beta < s,$ we instead write
    \begin{align*}
        I_{\Gamma_{\P}} &= n^{-r_\beta}\int_{\Gamma_{\P}}\frac{q(x)}{p(x)^{r_\beta}}\, dx \\
        &= n^{-r_\beta}\int_{\Gamma_{\P}}\Big(\frac{q(x)}{p(x)^{\gamma}}\Big)^aq(x)^{(1 - s)(1 - a)}p(x)^{-(r_\beta - \gamma a)}q(x)^{s(1 - a)}\, dx,
    \end{align*}
    with 
    \begin{align}
        \label{eq.a}
        a := \tfrac{r_\beta - s}{\gamma - s} \in (0, 1).    
    \end{align}
    We check that
    \begin{align*}
        r_\beta - \gamma a = \frac{r_\beta(\gamma - s) - \gamma(r_\beta -s)}{\gamma - s} = s\frac{\gamma - r_\beta}{\gamma - s}  = s(1 - a) > 0.
    \end{align*}
    Hence, using the definition of $\Gamma_{\P}$ before applying H\"older's inequality with conjugates $1/a$ and $1/(1 - a)$, we get
    \begin{align*}
        I_{\Gamma_{\P}} &\leq (\ell_{\P}C_{\P})^{\gamma a - r_\beta}(\ell_{\Q}C_q)^{s(1 - a)}n^{-\gamma a}m^{-s(1 - a)}\int_{\Gamma_{\P}}\Big(\frac{q(x)}{p(x)^{\gamma}}\Big)^aq(x)^{(1 - s)(1 - a)}\, dx\\
        &\leq (\ell_{\P}C_{p})^{\gamma a - r_\beta}(\ell_{\Q}C_q)^{s(1 - a)}n^{-\gamma a}m^{-s(1 - a)}\mT_{\P}(\gamma)^{a}\mT_{\Q}(s)^{1-a},
    \end{align*}
    which gives us the bound in the accelerated regime.\\
    
    \noindent\textbf{Bound on $\Gamma_{\Q}$.} This case is symmetric to the bound on $\Gamma_{\P}.$ Following the exact same steps and exchanging the roles of $\P_{\sX}$ and $\Q\mathstrut_{\!\sX}$ leads to the following. In the \textbf{wedge regime}, we obtain the bound
    \begin{align*}
        I_{\Gamma_{\Q}} &\leq \bigg[\frac{(C_q\ell_{\Q})^{r_T - r_\beta}}{m^{r_T}}\mT_{\Q}(r_T)\bigg] \wedge \bigg[(C_q\ell_{\Q})^{-r_\beta}\Big(\frac{C_p\ell_{\P}}{n}\Big)^{r_S}\mT_{\P}(r_S)\bigg],
    \end{align*}
    and, in the \textbf{accelerated regime}, we instead rewrite
    \begin{align*}
        I_{\Gamma_{\Q}} &= m^{-r_\beta}\int_{\Gamma_{\Q}}\frac{q(x)}{q(x)^{r\beta}}\, dx\\
        &= m^{-r_\beta}\int_{\Gamma_{\Q}}\Big(\frac{q(x)}{p(x)^\gamma}\Big)^aq(x)^{(1 - s)(1 - a)}p(x)^{\gamma a}q(x)^{-(r_\beta - s(1 - a))}\, dx.
    \end{align*}
    Then, we already know that $r_\beta - s(1 - a) = \gamma a >0,$ therefore, on $\Gamma_{\Q},$
    \begin{align*}
    q(x)^{-(r_\beta - s(1 - a))} \leq (m/(C_q\ell_{\Q}))^{r_\beta - s(1 - a)},
    \end{align*}
    and this enables to conclude
    \begin{align*}
        I_{\Gamma_{\Q}} &\leq (C_q\ell_{\Q})^{s(1 - a) - r_\beta}(C_p\ell_{\P})^{\gamma a}n^{-\gamma a}m^{-s(1 - a)}\mT_{\P}(\gamma)^{a}\mT_{\Q}(s)^{1 - a}.
    \end{align*}
    
    \noindent\textbf{Bound on $\Theta.$} We aim to bound
    \begin{align*}
        I_{\Theta} = \int_{\Theta}q(x)\, dx.
    \end{align*}
    We again distinguish the two regimes. In the \textbf{wedge regime}, using the fact that if $\gamma > 0$ then $\P_{\sX} \ll \Q\mathstrut_{\!\sX},$ it holds that 
    \begin{align*}
        I_{\Theta} = \int_{\Theta}q(x)\, dx &= \int_{\Theta}\Big(p(x)^{r_S}\frac{q(x)}{p(x)^{r_S}}\Big)\wedge \Big(q(x)^{r_T}\frac{q(x)}{q(x)^{r_T}}\Big)\, dx\\
        &\leq \bigg(\int_{\Theta}p(x)^{r_S}\frac{q(x)}{p(x)^{r_S}}\, dx\bigg)\wedge\bigg(q(x)^{r_T}\frac{q(x)}{q(x)^{r_T}}\, dx\bigg)\\
        &\leq \Big[\Big(\frac{C_p\ell_{\P}}{n}\Big)^{r_S}\mT_{\P}(r_S)\Big]\wedge \Big[\Big(\frac{C_q\ell_{\Q}}{m}\Big)^{r_T}\mT_{\Q}(r_T)\Big],
    \end{align*}
    where we also used that by definition, if $x \in \Theta,$ then $p(x) < C_p\ell_{\P}/n$ and $q(x) < C_q\ell_{\Q}/m.$ In the \textbf{accelerated regime}, we instead proceed as follows
    \begin{align*}
        I_{\Theta} = \int_{\Theta}q(x)\, dx &\leq \int_{\Theta} \Big(\frac{q(x)}{p(x)^{\gamma}}\Big)^a\Big(\frac{q(x)}{q(x)^s}\Big)^{1 - a}p(x)^{\gamma a}q(x)^{s(1 - a)}\, dx\\
        &\leq (C_p\ell_{\P})^{\gamma a}(C_q\ell_{\Q})^{s(1 - a)}n^{-\gamma a}m^{-s(1 - a)}\mT_{\P}(\gamma)^{a}\mT_{\Q}(s)^{1-a},
    \end{align*}
    while choosing $a$ as in \eqref{eq.a}.\\
    
    \noindent\textbf{Conclusion.} It remains to gather the terms to obtain, in the \textbf{wedge regime},
    \begin{align*}
        \big\|\wh f - f_*\big\|^2_{L^2(\Q\mathstrut_{\!\sX})} &\leq C_{\Delta}\log(nm)^{r_\beta}\bigg[\frac{1}{\k_{\P}n^{r_S}}(C_p\ell_{\P})^{r_S - r_\beta}\mT_{\P}(r_S)\bigg]\wedge\bigg[\frac{1}{\k_{\Q}m^{r_T}}(C_q\ell_{\Q})^{r_T - r_\beta}\mT_{\Q}(r_T)\bigg]\\
        &\qquad + C_{\Gamma_{\P}}\log(nm)^{r_\beta}\bigg[\frac{(C_p\ell_{\P})^{r_S - r_\beta}}{n^{r_S}}\mT_{\P}(r_S)\bigg] \wedge \bigg[(C_p\ell_{\P})^{-r_\beta}\Big(\frac{C_q\ell_{\Q}}{m}\Big)^{r_T}\mT_{\Q}(r_T)\bigg]\\
        &\qquad + C_{\Gamma_{\Q}}\log(nm)^{r_\beta}\bigg[(C_q\ell_{\Q})^{-r_\beta}\Big(\frac{C_p\ell_{\P}}{n}\Big)^{r_S}\mT_{\P}(r_S)\bigg]\wedge \bigg[\frac{(C_q\ell_{\Q})^{r_T - r_\beta}}{m^{r_T}}\mT_{\Q}(r_T)\bigg]\\
        &\qquad + C_{\Theta}\Big[\Big(\frac{C_p\ell_{\P}}{n}\Big)^{r_S}\mT_{\P}(r_S)\Big]\wedge \Big[\Big(\frac{C_q\ell_{\Q}}{m}\Big)^{r_T}\mT_{\Q}(r_T)\Big].
    \end{align*}
    Taking $\k_{\P} = \k_{\Q} = 1$ and recalling that $\ell_{\P} = \ell_{\Q} = \lceil V_\tau\log(nm)\rceil$ leads to first obtain $C_p = C_q =: C_0$ and
    \begin{align*}
        \big\|\wh f - f_*\big\|^2_{L^2(\Q\mathstrut_{\!\sX})} &\leq C_{\Delta}\bigg[(C_0V_\tau)^{r_S - r_\beta}\mT_{\P}(r_S)\Big(\frac{\log(nm)}{n}\Big)^{r_S}\bigg]\wedge\bigg[(C_0V_\tau)^{r_T - r_\beta}\mT_{\Q}(r_T)\Big(\frac{\log nm}{m}\Big)^{r_T}\bigg]\\
        &\qquad + C_{\Gamma_{\P}}\bigg[(C_0V_\tau)^{r_S - r_\beta}\mT_{\P}(r_S)\Big(\frac{\log (nm)}{n}\Big)^{r_S}\bigg] \wedge \bigg[(C_0V_\tau)^{r_T-r_\beta}\mT_{\Q}(r_T)\Big(\frac{\log(nm)}{m}\Big)^{r_T}\bigg]\\
        &\qquad + C_{\Gamma_{\Q}}\bigg[(C_0V_\tau)^{r_S-r_\beta}\mT_{\P}(r_S)\Big(\frac{\log(nm)}{n}\Big)^{r_S}\bigg]\wedge \bigg[(C_0V_\tau)^{r_T - r_\beta}\mT_{\Q}(r_T)\Big(\frac{\log(nm)}{m}\Big)^{r_T}\bigg]\\
        &\qquad + C_{\Theta}\Big[(C_0V_\tau)^{r_S}\mT_{\P}(r_S)\Big(\frac{\log(nm)}{n}\Big)^{r_S}\Big]\wedge \Big[(C_0V_\tau)^{r_T}\mT_{\Q}(r_T)\Big(\frac{\log(nm)}{m}\Big)^{r_T}\Big]\\
        &\leq \ol C_0\Big[\mT_{\P}(r_S)\Big(\frac{\log(nm)}{n}\Big)^{r_S}\Big]\wedge \Big[\mT_{\Q}(r_T)\Big(\frac{\log(nm)}{m}\Big)^{r_T}\Big],
    \end{align*}
    for $\ol C_0$ being the sum of all the constants above.\\
    \noindent In the \textbf{accelerated regime}, collecting the terms leads to
    \begin{align*}
        \big\|\wh f - f_*\big\|_{L^2(\Q\mathstrut_{\!\sX})}^2 &\leq C_{\Delta}\log(nm)^{r_\beta}\k_{\P}^{-\gamma a/r_{\beta}}\k_{\Q}^{-s(1 - a)/r_\beta}n^{-\gamma a}m^{-s(1- a)}\mT_{\P}(\gamma)^{a} \mT_{\Q}(s)^{1 - a}\\
        &\quad + C_{\Gamma_{\P}}\log(nm)^{r_\beta}(\ell_{\P}C_{p})^{\gamma a - r_\beta}(\ell_{\Q}C_q)^{s(1 - a)}n^{-\gamma a}m^{-s(1 - a)}\mT_{\P}(\gamma)^{a}\mT_{\Q}(s)^{1-a}\\
        &\quad + C_{\Gamma_{\Q}}\log(nm)^{r_\beta}(C_q\ell_{\Q})^{s(1 - a) - r_\beta}(C_p\ell_{\P})^{\gamma a}n^{-\gamma a}m^{-s(1 - a)}\mT_{\P}(\gamma)^{a}\mT_{\Q}(s)^{1 - a}\\
        &\quad + C_{\Theta}(C_p\ell_{\P})^{\gamma a}(C_q\ell_{\Q})^{s(1 - a)}n^{-\gamma a}m^{-s(1 - a)}\mT_{\P}(\gamma)^{a}\mT_{\Q}(s)^{1-a}.
    \end{align*}
    Picking $\k_{\P} = \k_{\Q} = 1$ and recalling that $\ell_{\P} = \ell_{\Q} = \lceil V_\tau \log(nm)\rceil$ leads to $C_p = C_q$ and
    \begin{align*}
        \big\|\wh f - f_*\big\|_{L^2(\Q\mathstrut_{\!\sX})}^2 &\leq C_{\Delta}\Big(\frac{\log(nm)}{n}\Big)^{\gamma a}\Big(\frac{\log(nm)}{m}\Big)^{s(1- a)}\mT_{\P}(\gamma)^{a} \mT_{\Q}(s)^{1 - a}\\
        &\quad + C_{\Gamma_{\P}}\Big(\frac{\log(nm)}{n}\Big)^{\gamma a}\Big(\frac{\log(nm)}{m}\Big)^{s(1- a)}\mT_{\P}(\gamma)^{a}\mT_{\Q}(s)^{1-a}\\
        &\quad + C_{\Gamma_{\Q}}\Big(\frac{\log(nm)}{n}\Big)^{\gamma a}\Big(\frac{\log(nm)}{m}\Big)^{s(1- a)}\mT_{\P}(\gamma)^{a}\mT_{\Q}(s)^{1 - a}\\
        &\quad + C_{\Theta}(C_0V_\tau)^{r_\beta}\Big(\frac{\log(nm)}{n}\Big)^{\gamma a}\Big(\frac{\log(nm)}{m}\Big)^{s(1- a)}\mT_{\P}(\gamma)^{a}\mT_{\Q}(s)^{1-a}\\
        &= \tilde C_0\Big(\frac{\log(nm)}{n}\Big)^{\gamma a}\Big(\frac{\log(nm)}{m}\Big)^{s(1- a)}\mT_{\P}(\gamma)^{a}\mT_{\Q}(s)^{1-a},
    \end{align*}
    which yields the claimed upper bound after replacing $a = \tfrac{r_\beta - s}{\gamma - s}.$ Recall that we were working in the event $E_n^{\P_{\sX}}(\ell_{\P}, \wh f_{\P}) \cap E_m^{\Q\mathstrut_{\!\sX}}(\ell_{\Q}, \wh f_{\Q}).$ It therefore remains to use Propositions \ref{prop.neighbours.distance.bound.zeta} and \ref{prop.neighbours.distance.lower.bound.zeta}, Corollary \ref{cor.variance.bound.nice.version}, and Remark \ref{rem.bound.variable.k} together with the union bound to obtain
    \begin{align*}
        \PP\bigg\{E_n^{\P_{\sX}}(\ell_{\P}, \wh f_{\P}) \cap E_m^{\Q\mathstrut_{\!\sX}}(\ell_{\Q}, \wh f_{\Q})\bigg\} \geq 1 - 2n^{-\tau} - n^{1-\tau} - 2m^{-\tau} - m^{1-\tau} \geq 1 - 3n^{1-\tau} - 3m^{1-\tau}.
    \end{align*}
\end{proof}

\end{document}
